\chapter{Język i podstawowe konstrukcje topologiczne}
\label{ch:wstep-topologia}

Topologia algebraiczna bada przestrzenie za pomocą niezmienników, które można
obliczać algebraicznie: grup podstawowych, grup homologii i pierścieni
kohomologii. Nie mierzymy odległości; pytamy, jak przestrzeń jest sklejona.
Celem tego rozdziału jest zebranie podstawowego języka potrzebnego w dalszej
części książki. Zakładamy elementarną znajomość zbiorów, odwzorowań i algebry
liniowej. Wszystkie konstrukcje poniżej dotyczą przestrzeni topologicznych.

\section{Przestrzenie topologiczne, ciągłość i przykłady}

\begin{definicja}
Topologią na zbiorze $X$ nazywamy rodzinę $\tau\subseteq\mathcal P(X)$
zawierającą $\varnothing$ i $X$, zamkniętą na dowolne sumy oraz skończone
przecięcia. Parę $(X,\tau)$ nazywamy przestrzenią topologiczną; elementy
$\tau$ nazywamy zbiorami otwartymi. Odwzorowanie $f:X\to Y$ jest ciągłe,
jeśli przeciwobraz każdego otwartego podzbioru $Y$ jest otwarty w $X$.
Homeomorfizmem nazywamy ciągłą bijekcję o ciągłej odwrotności.
\end{definicja}

\begin{przyklad}
Podstawowe modele to $\R^n$ z topologią euklidesową, sfera
$S^n=\{x\in\R^{n+1}:\|x\|=1\}$ i dysk $D^n=\{x\in\R^n:\|x\|\leq1\}$.
Nie utożsamiamy $S^0$ z okręgiem: $S^0$ składa się z dwóch punktów.
Punkty $S^1$ można zapisać jako $e^{2\pi it}$ dla $t\in\R$.
\end{przyklad}

\begin{definicja}
Podprzestrzeń $A\subseteq X$ wyposażamy w topologię $\{A\cap U:U\in\tau_X\}$.
Na iloczynie $X\times Y$ przyjmujemy topologię generowaną przez prostokąty
$U\times V$ z $U\subseteq X$, $V\subseteq Y$ otwartymi. Przestrzeń $X$
nazywamy spójną, gdy nie jest sumą dwóch rozłącznych, niepustych zbiorów
otwartych; nazywamy ją łukowo spójną, gdy każde dwa punkty można połączyć
ciągłą drogą $[0,1]\to X$.
\end{definicja}

\begin{stwierdzenie}
Niech $X$ będzie łukowo spójna. Wówczas $X$ jest spójna.
\end{stwierdzenie}
\begin{proof}
Przypuśćmy, że $X=U\sqcup V$, gdzie $U,V$ są niepuste i otwarte.
Wybierzmy $u\in U$ i $v\in V$ oraz drogę $\gamma$ od $u$ do $v$.
Zbiory $\gamma^{-1}(U)$ oraz $\gamma^{-1}(V)$ byłyby rozłącznymi,
niepustymi zbiorami otwartymi o sumie $[0,1]$. Przeczyłoby to spójności
odcinka $[0,1]$.
\end{proof}

\section{Topologia ilorazowa i przestrzenie sklejone}

\begin{definicja}[Topologia ilorazowa]
\label{def:topologia-ilorazowa}
Niech $\sim$ będzie relacją równoważności na $X$ i
$q:X\to X/{\sim}$ odwzorowaniem $x\mapsto[x]$.
Zbiór $U\subseteq X/{\sim}$ nazywamy otwartym wtedy i tylko wtedy, gdy
$q^{-1}(U)$ jest otwarty w $X$. Tak otrzymaną topologię nazywamy
\emph{topologią ilorazową}.
\end{definicja}

\begin{twierdzenie}[Własność uniwersalna ilorazu]
Niech $q:X\to X/{\sim}$ będzie odwzorowaniem ilorazowym i niech $Y$ będzie
przestrzenią topologiczną. Odwzorowanie $F:X/{\sim}\to Y$ jest ciągłe wtedy
i tylko wtedy, gdy $F\circ q:X\to Y$ jest ciągłe.
W szczególności każde ciągłe $f:X\to Y$, stałe na klasach relacji $\sim$,
wyznacza jedyne ciągłe $\bar f:X/{\sim}\to Y$ takie, że $f=\bar f\circ q$.
\end{twierdzenie}
\begin{proof}
Gdy $F$ jest ciągłe, ciągłe jest także $F\circ q$. W drugą stronę,
dla każdego otwartego $V\subseteq Y$ mamy
$q^{-1}(F^{-1}(V))=(F\circ q)^{-1}(V)$, co jest otwarte.
Z definicji topologii ilorazowej wynika otwartość $F^{-1}(V)$.
Jeśli $f$ jest stałe na klasach, wzór $\bar f([x])=f(x)$ nie zależy od wyboru
przedstawiciela; ciągłość wynika z pierwszej części dowodu.
\end{proof}

\begin{przyklad}[Okrąg i torus]
Iloraz $[0,1]/(0\sim1)$ jest homeomorficzny z $S^1$ przez
$[t]\mapsto e^{2\pi it}$; ciągła bijekcja jest homeomorfizmem, ponieważ
jej dziedzina jest zwarta, a przeciwdziedzina Hausdorffa.
Podobnie $\R^2/\Z^2$, gdzie $(x,y)\sim(x+m,y+n)$ dla $m,n\in\Z$,
jest torusem $T^2=S^1\times S^1$.
\end{przyklad}

\begin{definicja}
Dla pary $(X,A)$ z $A\subseteq X$ przestrzeń $X/A$ powstaje przez
utożsamienie wszystkich punktów $A$ w jeden punkt. Dla przestrzeni z
punktami bazowymi $(X,x_0)$ i $(Y,y_0)$ \emph{suma klinowa}
$X\vee Y$ jest ilorazem $X\sqcup Y$ przez relację $x_0\sim y_0$.
Stożek, zawieszenie i produkt smash dla dobrze bazowanych
przestrzeni zapisujemy odpowiednio jako
$CX=X\times[0,1]/(X\times\{1\})$,
$\Sigma X=(X\times[0,1])/(X\times\{0\},X\times\{1\})$,
oraz $X\wedge Y=(X\times Y)/(X\vee Y)$,
gdzie $X\vee Y$ utożsamiamy z
$(X\times\{y_0\})\cup(\{x_0\}\times Y)$.
W zawieszeniu obie końcowe kopie $X$ są sklejane do \emph{dwóch różnych}
punktów.
\end{definicja}

\begin{uwaga}
Powyższy stożek $CX$ jest stożkiem niezredukowanym (jedna sklejona
podstawa). Termin „zredukowany” jest zarezerwowany dla dodatkowego
utożsamienia $\{x_0\}\times[0,1]$; w dalszej części nie będziemy
mylić tych dwóch konstrukcji.
\end{uwaga}

\begin{przyklad}
Zawieszenie $S^n$ jest homeomorficzne z $S^{n+1}$, a zawieszenie $S^0$
z dwiema rozłącznymi kopiami punktów daje okrąg $S^1$.
\end{przyklad}

\section{Homotopie i deformacyjne retrakcje}

\begin{definicja}
Dwa ciągłe odwzorowania $f,g:X\to Y$ są homotopijne, piszemy
$f\simeq g$, jeżeli istnieje ciągłe $H:X\times[0,1]\to Y$ takie, że
$H(x,0)=f(x)$ i $H(x,1)=g(x)$ dla każdego $x\in X$.
Podzbiór $A\subseteq X$ jest \emph{deformacyjnym retraktem} $X$, jeśli
istnieje homotopia $H:\id_X\simeq i\circ r$, gdzie $r:X\to A$ jest
retrakcją, $i:A\hookrightarrow X$ jest inkluzją oraz
$H(a,t)=a$ dla wszystkich $a\in A$, $t\in[0,1]$.
Przestrzeń jest ściągalna, gdy jej identyczność jest homotopijna z
odwzorowaniem stałym.
\end{definicja}

\begin{przyklad}[Dziurawa płaszczyzna]
Płaszczyzna $\R^2\setminus\{0\}$ ma deformacyjny retrakt $S^1$.
Dla $x\ne0$ kładziemy
$H(x,t)=\big((1-t)+t/\|x\|\big)x$.
Współczynnik jest dodatni, więc homotopia nie przechodzi przez zero;
$H(x,0)=x$, $H(x,1)=x/\|x\|$ i $H(a,t)=a$ dla $a\in S^1$.
Ten rachunek tłumaczy, dlaczego oba zbiory mają takie same grupy homotopii
i homologii, mimo że nie są homeomorficzne.
\end{przyklad}

\begin{twierdzenie}
Jeżeli $A$ jest deformacyjnym retraktem $X$, to inkluzja $i:A\to X$
jest równoważnością homotopijną, czyli ma homotopijną odwrotność.
\end{twierdzenie}
\begin{proof}
Retrakcja $r:X\to A$ spełnia $r\circ i=\id_A$.
Z definicji deformacyjnej retrakcji mamy $i\circ r\simeq\id_X$.
To właśnie dwa warunki definiujące odwrotność homotopijną.
\end{proof}

\section{Przestrzenie rzutowe i modele do obliczeń}

\begin{definicja}
Rzeczywista przestrzeń rzutowa $\R P^n$ to iloraz $S^n$ przez
utożsamienie punktów antypodycznych $x\sim-x$.
Zespolona przestrzeń rzutowa $\C P^n$ to zbiór zespolonych prostych
w $\C^{n+1}$ z topologią ilorazową pochodzącą z
$\C^{n+1}\setminus\{0\}$ i relacji $z\sim\lambda z$ dla $\lambda\in\C^*$.
\end{definicja}

\begin{uwaga}
Przestrzenie rzutowe i torusy będą naszymi laboratoriami: najpierw
rozpoznamy ich rozkłady komórkowe, następnie policzymy homologię oraz
pierścienie kohomologii. Ważne będzie odróżnienie niezmienników
\emph{grupowych} od dodatkowych struktur, takich jak iloczyn kubkowy.
\end{uwaga}

% Rozdzial przygotowany do samodzielnego skryptu topologii algebraicznej.
